\chapter{Données et feature stores}\label{ch:ml:data-and-feature-stores}

Une prévision du carnet d'ordres à horizon d'une seconde affiche un coefficient d'information de rang (rank IC) de 0.44 dans son backtest et de 0.36 en production. Rien n'est
faux dans le modèle. En recherche, ses variables ont été calculées sur des barres d'une seconde et rattachées à chaque décision à la fin de
sa seconde~; en production, elles sont calculées à partir de chaque message jusqu'à la décision. Les variables de recherche avaient vu jusqu'à une
seconde du futur, et le modèle avait appris à l'utiliser. Ce chapitre vise à rendre cela impossible~: déclarer chaque
variable une seule fois et la calculer de la même manière pour l'entraînement et pour le trading, la stocker avec son temps de connaissance, connaître
la fraîcheur de chaque valeur, et tester en continu que les deux calculs concordent. Ses mesures utilisent le marché synthétique à carnet d'ordres du livre et
quatre écarts injectés~; le test de parité les détecte tous, et un seul d'entre eux coûte quelque chose au modèle, ce qui est
une raison de les détecter tous plutôt que de deviner lesquels comptent.

\section{La même variable deux fois~: hors ligne et en ligne}

\begin{definition}[Feature store, magasin hors ligne, magasin en ligne, définition de variable]\label{def:ml:data-and-feature-stores:store}
Un \emph{feature store}\index{feature store} est le système qui calcule, stocke et sert les entrées des modèles d'une entreprise à partir d'un seul
jeu de déclarations. Son \emph{magasin hors ligne}\index{magasin hors ligne} (offline store) conserve les valeurs des variables sur l'historique, indexées par entité et
instant, pour l'entraînement et les backtests~; son \emph{magasin en ligne}\index{magasin en ligne} (online store) conserve les dernières valeurs, mises à jour à mesure que les données
arrivent, pour les modèles en production. Une \emph{définition de variable}\index{définition de variable} (feature definition) déclare une variable sous forme de donnée (ses
entrées, son agrégation, sa fenêtre, son horloge et sa version) afin que les deux magasins la calculent à partir de la même déclaration
(Sculley et coauteurs, 2015, sur les raisons de cet enjeu~; Hermann et Del Balso, 2017, pour un des premiers dispositifs de production).
\end{definition}

Le chapitre déclare douze variables du carnet d'ordres (\cref{tab:ml:fs:skew})~: la fourchette, les déséquilibres de file et de profondeur
et le mid pondéré (les dernières valeurs après chaque message, issues du moteur en flux \texttt{firm.lobfeat} du livre~7), le déséquilibre du flux
d'ordres sur cinq et trente secondes, les volumes échangés et signés sur trente secondes, les nombres de transactions et de messages sur dix secondes, et
la variation du mid sur cinq et trente secondes. Chacune est un \texttt{FeatureDef}~: un nom, l'entrée par événement, une agrégation (dernière
valeur, somme, nombre, variation) et une fenêtre. Le \omterm{def:ml:data-and-feature-stores:store}{magasin hors ligne} calcule les douze pour un historique entier, sur une liste d'instants de décision
en une seule fois, à partir de sommes cumulées et de jointures as-of (\cref{lst:ml:fs:offline})~; le moteur en ligne consomme les messages un à un
et tient à jour des fenêtres glissantes (\cref{lst:ml:fs:online}). Sur huit séances de trente minutes de \texttt{firm.tape}, 313\,681 messages
et 7\,714 décisions (une à chaque transaction), les deux concordent exactement à chaque décision~: tel est le contrat du magasin.

\section{Exactitude à date dans le magasin}

Chaque valeur du \omterm{def:ml:data-and-feature-stores:store}{magasin hors ligne} porte l'instant où elle est devenue connue, et une requête pour une décision à l'instant $t$ ne voit que les valeurs
connues à $t$~: la jointure as-of du livre~7 (chapitre~3), appliquée aux variables. Pour les variables pilotées par événements, le temps de connaissance est
celui du dernier message lu par la variable~; pour les variables issues de données plus lentes (fondamentaux, données alternatives, chapitre~15)
c'est l'instant de livraison, et le magasin est bitemporel. Les deux erreurs que cela écarte sont celles qui sont injectées ici~: une chaîne de recherche
qui lit un message au-delà de la décision (un décalage d'un cran dans une jointure), et une autre qui calcule les variables sur une horloge plus grossière
horloge et rattache à chaque décision la valeur en fin de sa barre.

\section{Fraîcheur et matérialisation}

\begin{definition}[Matérialisation, fraîcheur des variables]\label{def:ml:data-and-feature-stores:fresh}
La \emph{matérialisation}\index{matérialisation} est le traitement qui calcule les valeurs des variables à partir de leurs définitions et les écrit
dans un magasin, par lots pour le \omterm{def:ml:data-and-feature-stores:store}{magasin hors ligne} ou en continu pour le \omterm{def:ml:data-and-feature-stores:store}{magasin en ligne}.
La \emph{fraîcheur des variables}\index{fraîcheur des variables} (feature freshness) est l'âge d'une valeur servie~: le temps écoulé depuis la plus récente de ses entrées, grandeur qu'un système de production
surveille et borne.
\end{definition}

La fraîcheur est ce qui distingue la production de la recherche, même lorsque le code est partagé~: un flux d'entrée qui arrive en retard rend
périmée toute valeur en ligne qui en dépend. Le quatrième écart injecté retarde les publications de transactions de 200~millisecondes par rapport au
carnet, comme le ferait un flux de transactions distinct~; les trois variables de nombre de transactions et de volume diffèrent alors de la recherche dans 92 à 94~\% des
décisions, de 7 à 17~\% de leur écart-type, alors que les variables de carnet ne sont pas touchées.

\section{Détection des écarts}

\begin{definition}[Écart entraînement--service]\label{def:ml:data-and-feature-stores:skew}
Un \emph{écart entraînement--service}\index{ecart@écart entraînement--service} (training--serving skew) est toute différence entre les valeurs des variables sur lesquelles un modèle a été entraîné
et celles qu'il reçoit en production au même instant~: code différent, horloges différentes, entrées différentes,
arithmétique différente ou synchronisation des données différente.
\end{definition}

Le test de parité (\cref{lst:ml:fs:parity}) recalcule hors ligne les valeurs de la production pour un échantillon de décisions récentes et les compare
variable par variable. Le \Cref{tab:ml:fs:skew} montre l'empreinte de chaque écart injecté, et le \cref{tab:ml:fs:parity} ce que voient le test
et le modèle. Un contrôle de cinq décisions tirées au hasard par séance, exécuté cinq fois sur chacune des huit séances, détecte
tous les écarts, à chaque fois~: chacun modifie la plupart des décisions. La tolérance est un choix de conception~: à $10^{-9}$, le test signale aussi
le stockage float32 de la production (une variation d'au plus $10^{-7}$, sans effet sur le modèle), et à $10^{-6}$ il cesse de le signaler tout en
continuant à signaler les trois autres dans 96 à 100~\% des décisions.

\begin{table}[htbp]
\centering\small
\setlength{\tabcolsep}{4pt}%
\begin{tabular}{@{}>{\raggedright\arraybackslash}p{4.2cm}>{\centering\arraybackslash}p{2.5cm}>{\centering\arraybackslash}p{2.5cm}>{\centering\arraybackslash}p{2.5cm}>{\centering\arraybackslash}p{2.5cm}@{}}
\toprule
& \multicolumn{4}{c}{Écart quadratique moyen en proportion de l'écart-type de la variable}\\
\cmidrule(l){2-5}
variable & un événement d'avance & barres d'une seconde & float32 & transactions en retard\\
\midrule
fourchette & 0.38 & 0.92 & 0 & 0\\
déséquilibre de file & 0.24 & 0.73 & 0.00 & 0\\
déséquilibre de profondeur & 0.11 & 0.41 & 0.00 & 0\\
mid pondéré moins mid & 0.30 & 0.79 & 0.00 & 0\\
OFI 5 s & 0.03 & 0.24 & 0 & 0\\
OFI 30 s & 0.01 & 0.07 & 0 & 0\\
volume 30 s & 0.01 & 0.07 & 0 & 0.07\\
volume signé 30 s & 0.02 & 0.10 & 0 & 0.11\\
transactions 10 s & 0.02 & 0.18 & 0 & 0.17\\
variation du mid 5 s & 0.09 & 0.39 & 0 & 0\\
variation du mid 30 s & 0.03 & 0.13 & 0 & 0\\
messages 10 s & 0.00 & 0.04 & 0 & 0\\
\bottomrule
\end{tabular}
\caption{Ampleur de chaque écart injecté, variable par variable, sur 7\,714 décisions (0.00~: une différence inférieure à
un demi-centième~; 0~: aucune). Données~: \texttt{ml\_featstore.feature\_skew}.}\label{tab:ml:fs:skew}
\end{table}

\begin{table}[htbp]
\centering\small
\setlength{\tabcolsep}{4pt}%
\begin{tabular}{@{}>{\raggedright\arraybackslash}p{4.8cm}>{\centering\arraybackslash}p{1.9cm}>{\centering\arraybackslash}p{2.4cm}>{\centering\arraybackslash}p{1.6cm}>{\centering\arraybackslash}p{1.7cm}>{\centering\arraybackslash}p{2.2cm}@{}}
\toprule
& décisions & détectées & \multicolumn{2}{c}{rank IC} & P\&L en production\\
\cmidrule(lr){4-5}
écart & discordantes & (contrôle sur 5 échantillons) & backtest & production & (ticks)\\
\midrule
aucun & 0 & 0 & 0.411 & 0.411 & 0.089\\
recherche, un événement d'avance & 100~\% & 100~\% & 0.395 & 0.418 & 0.091\\
recherche sur barres d'une seconde & 96~\% & 100~\% & 0.438 & 0.360 & 0.078\\
stockage float32 en production & 92~\% & 100~\% & 0.411 & 0.411 & 0.089\\
production, publications de transactions avec 200 ms de retard & 100~\% & 100~\% & 0.411 & 0.411 & 0.089\\
\bottomrule
\end{tabular}
\caption{Test de parité (tolérance $10^{-9}$) et prévision ridge de la variation du mid sur la seconde suivante, entraînée sur quatre
séances avec chaque chaîne de recherche et testée sur quatre autres~: son backtest sur les variables de recherche et son résultat en production sur
variables de production, rank IC et P\&L moyen par décision de trading du signe de la prévision. Données~:
\texttt{ml\_featstore.parity\_table}, \texttt{model\_effects}.}\label{tab:ml:fs:parity}
\end{table}

\begin{omfigure}
\begin{tikzpicture}
\begin{axis}[width=11.5cm,height=5.8cm,ybar,bar width=10pt,ymin=0.3,ymax=0.46,ylabel={rank IC, prévision à une seconde},area legend,
  xtick={0,1,2,3,4},xticklabels={{aucun},{un événement\\ en avance},{barres\\ d'une seconde},{stockage\\ float32},{transactions\\ en retard}},
  xticklabel style={align=center},tick label style={font=\small},label style={font=\small},xmin=-0.5,xmax=4.5,grid=major,
  grid style={gray!25},legend columns=2,legend style={font=\footnotesize,at={(0.5,-0.32)},anchor=north,draw=none,fill=none,/tikz/every even column/.append style={column sep=8pt}},
  table/col sep=comma]
\addplot[fill=omThm,draw=none,bar shift=-5.5pt] table[x=i,y=backtest]{figdata/ml/24-data-and-feature-stores/ic.csv};
\addplot[fill=omDef,draw=none,bar shift=5.5pt] table[x=i,y=live]{figdata/ml/24-data-and-feature-stores/ic.csv};
\legend{backtest (variables de recherche),en direct (variables de production)}
\end{axis}
\end{tikzpicture}

\omcaption{Ce que chaque écart fait au modèle~: le rank IC que son backtest de recherche annonce et celui qu'il obtient en direct. Données~:
\texttt{ml\_featstore.model\_effects}.}\label{fig:ml:fs:ic}
\end{omfigure}

Le résultat qui donne son nom au chapitre figure à la troisième ligne du tableau. La recherche sur barres d'une seconde gonfle l'IC du backtest de 0.411 à
0.438 et, parce que le modèle a appris à s'appuyer sur des variables ayant vu le futur, fournit 0.360 en direct, en deçà de ce que gagne la
chaîne propre~; le P\&L par décision passe de 0.089 tick à 0.078. L'anticipation d'un seul événement produit ici l'effet inverse~: le
message qui suit une transaction est généralement une mise à jour de cotation au fond du carnet, qui ajoute du bruit plutôt que de l'information, et le backtest
sous-estime le résultat en direct (0.395 contre 0.418). Sur la prévision à cinq secondes, les mêmes écarts déplacent l'IC de moins de 0.01~: 
qu'un écart coûte quelque chose dépend de la corrélation de son erreur avec la cible, ce qui n'est pas connaissable à l'avance. Le
rôle du magasin est de faire concorder les magasins, non de juger quels désaccords sont inoffensifs.

\begin{method}[Calcul des variables en production]\label{met:ml:fs:method}
\begin{enumerate}
\item Déclarer chaque variable une seule fois, sous forme de données, avec son horloge et sa fenêtre, et la versionner.
\item Calculer hors ligne et en ligne à partir des mêmes déclarations~; dans le \omterm{def:ml:data-and-feature-stores:store}{magasin hors ligne}, horodater chaque valeur de son temps de connaissance
et de son point de jointure dans le temps.
\item Lancer chaque jour un test de parité sur un échantillon de décisions récentes, avec des tolérances fixées par variable, et alerter à la moindre divergence.
\item Surveiller la fraîcheur de chaque variable en ligne par rapport à une borne, et replier ou arrêter lorsque les entrées deviennent périmées.
\end{enumerate}
\end{method}

\section{Tutoriel~: les mêmes variables deux fois}\label{tut:ml:data-and-feature-stores:tut}

\begin{tutorial}
\textbf{Objectif.} Déclarer douze variables, les calculer hors ligne et en ligne sur huit séances, introduire quatre écarts, et mesurer ce que
voient le test de parité et le modèle. \textbf{État final~:} \cref{tab:ml:fs:skew,tab:ml:fs:parity}, \cref{fig:ml:fs:ic}.
\begin{tutsteps}
\item \textbf{Le \omterm{def:ml:data-and-feature-stores:store}{magasin hors ligne}~: toutes les décisions à la fois.}
\omcode{code/firm/featstore/firm_featstore.py}{61}{79}{Valeurs vectorisées des variables au point dans le temps.}\label{lst:ml:fs:offline}
\item \textbf{Le moteur en ligne~: fenêtres glissantes.}
\omcode{code/firm/featstore/firm_featstore.py}{108}{124}{Valeurs en ligne à tout instant.}\label{lst:ml:fs:online}
\item \textbf{Le test de parité.}
\omcode{code/firm/featstore/firm_featstore.py}{143}{151}{Comparaison des magasins sur un échantillon de décisions.}\label{lst:ml:fs:parity}
\item \textbf{Exécuter} \texttt{ml\_featstore.parity\_table()}, \texttt{parity\_table(tol=1e-6)}, \texttt{model\_effects()},
\texttt{model\_effects(5.0)}, \texttt{feature\_skew()} et \texttt{fig\_featstore.py}.
\end{tutsteps}
\textbf{Pour aller plus loin.} Introduire un écart qui n'affecte qu'une décision sur mille et déterminer la taille d'échantillon qu'un contrôle quotidien doit avoir
pour le détecter en une semaine.
\end{tutorial}

\section{Construction~: le feature store}\label{bld:ml:data-and-feature-stores:featstore}

\begin{build}
\bfield{Rôle} Entrées du modèle calculées une seule fois selon leur définition, au point dans le temps hors ligne, fraîches en ligne, et vérifiées en continu
l'une par rapport à l'autre.
\bfield{Interface} \texttt{FeatureDef(name, input, agg, window, version)}, \texttt{events\_from\_tape(tape, levels)} (dans
\texttt{firm.lobfeat}), \texttt{offline(events, defs, times, lag\_events)}, \texttt{OnlineEngine(defs)} avec \texttt{on} et
\texttt{values}, \texttt{stream(events, defs, times)}, \texttt{parity(off, on, tol, sample, seed)},
\texttt{freshness(engine, now)}.
\bfield{Règles} Aucune variable n'existe en dehors d'une définition~; le \omterm{def:ml:data-and-feature-stores:store}{magasin hors ligne} ne renvoie jamais une valeur connue après l'instant de la requête~;
le test de parité tourne chaque jour~; les entrées périmées sont des alarmes.
\bfield{Tests d'acceptation} \texttt{code/firm/featstore/tests/}~: hors ligne et en ligne égaux exactement sur une séance~; chaque agrégation sur fenêtre
coïncide avec un recalcul par force brute~; une anticipation d'un seul événement est signalée~; la fraîcheur croît lorsque les entrées cessent.
\bfield{Pour aller plus loin} Un \omterm{def:ml:data-and-feature-stores:store}{magasin hors ligne} bitemporel sur \texttt{firm.pit} pour les variables lentes~; des tolérances par variable apprises de
l'arithmétique~; des moteurs en ligne en C++ ou en Rust, vérifiés contre la référence Python, comme le livre~7 l'a fait pour \texttt{firm.lobfeat}.
\end{build}

\begin{omsources}
\item D.~Sculley et coauteurs, ``Hidden technical debt in machine learning systems'', \emph{Advances in Neural Information
Processing Systems} 28, 2015.
\item J.~Hermann et M.~Del Balso, ``Meet Michelangelo: Uber's machine learning platform'', blog d'ingénierie d'Uber, 2017.
\item N.~Polyzotis, S.~Roy, S.~E.~Whang et M.~Zinkevich, ``Data lifecycle challenges in production machine learning: a
survey'', \emph{SIGMOD Record} 47(2), 2018.
\end{omsources}

\section{Exercices}

\begin{exercise}[$\star$]\label{exo:ml:data-and-feature-stores:1}
Une variable de volume sur trente secondes est calculée à une décision à 10:00:07.300. Quelles transactions le \omterm{def:ml:data-and-feature-stores:store}{magasin en ligne} inclut-il, et lesquelles
un pipeline de recherche sur barres d'une seconde, joint à la fin de la seconde de la décision, inclut-il~?
\end{exercise}

\begin{exercise}[$\star$]\label{exo:ml:data-and-feature-stores:2}
Un écart modifie une décision sur mille. Combien de décisions tirées au hasard un contrôle doit-il comparer pour le détecter avec
une probabilité de 95~\%~?
\end{exercise}

\begin{exercise}[$\star$]\label{exo:ml:data-and-feature-stores:3}
Pourquoi le stockage en float32 modifie-t-il les variables de déséquilibre mais pas les variables de volume et de nombre~?
\end{exercise}

\begin{exercise}[$\star\star$]\label{exo:ml:data-and-feature-stores:4}
Pourquoi les barres d'une seconde ont-elles nui à la prévision à une seconde et non à celle à cinq secondes~?
\end{exercise}

\begin{exercise}[$\star\star$]\label{exo:ml:data-and-feature-stores:5}
Pourquoi l'anticipation d'un seul événement a-t-elle dégradé le backtest au lieu de l'améliorer~? Le ferait-elle toujours~?
\end{exercise}

\begin{exercise}[$\star\star$]\label{exo:ml:data-and-feature-stores:6}
\emph{Trouver la faille.} «~Nous avons vérifié que les variables de production et de recherche ont la même moyenne et le même écart-type chaque jour,
donc il n'y a pas d'\omterm{def:ml:data-and-feature-stores:skew}{écart entraînement--service}.~»
\end{exercise}

\begin{exercise}[$\star\star\star$]\label{exo:ml:data-and-feature-stores:7}
\emph{Programmation.} Réentraîner le modèle à une seconde sur les valeurs que le moteur de production calcule lorsqu'il rejoue les séances d'entraînement
au lieu des variables de recherche sur barres d'une seconde, puis l'évaluer en direct. Que récupère-t-il, et que coûte-t-il~?
\end{exercise}

\begin{exercise}[$\star\star\star$]\label{exo:ml:data-and-feature-stores:8}
Concevoir la tolérance du test de parité pour une somme sur une fenêtre d'entrées float64 dont les calculs en ligne et hors ligne
additionnent les termes dans des ordres différents.
\end{exercise}

\section{Problème~: Deux fois les mêmes variables}

\begin{problem}[{Problème du week-end --- une définition, deux calculs}]\label{pb:ml:data-and-feature-stores:1}
Les douze variables du chapitre, huit séances et quatre écarts.

\medskip\noindent\textbf{Partie I --- Le magasin.}
\begin{enumerate}
\item Que contient une \omterm{def:ml:data-and-feature-stores:store}{définition de variable}, et pourquoi la traiter comme des données~?
\item En quoi les calculs hors ligne et en ligne diffèrent-ils, et comment savoir qu'ils concordent~?
\item Quel est le temps de connaissance d'une variable de carnet d'ordres, et d'un fondamental trimestriel~?
\item Qu'est-ce que la fraîcheur, et qu'est-ce qui l'a fait échouer ici~?
\end{enumerate}

\medskip\noindent\textbf{Partie II --- Les écarts.}
\begin{enumerate}[resume]
\item Décrire les quatre écarts et les variables que chacun touche.
\item Lesquels sont des bogues de recherche et lesquels des bogues de production~?
\item Quelle part des décisions chacun modifie-t-il~?
\item Comment la tolérance du test de parité modifie-t-elle ce qu'il signale~?
\end{enumerate}

\medskip\noindent\textbf{Partie III --- Le modèle.}
\begin{enumerate}[resume]
\item Quels sont les IC du backtest et en direct sous chaque écart~?
\item Pourquoi l'écart lié aux barres d'une seconde gonfle-t-il le backtest et dégrade-t-il le résultat en direct~?
\item Que change-t-il à un horizon de cinq secondes~?
\item À quoi ressemblerait l'écart d'IC sur le tableau de bord d'un desk~?
\end{enumerate}

\medskip\noindent\textbf{Partie IV --- Le verdict.}
\begin{enumerate}[resume]
\item Énoncer le \emph{résultat nommé}~: l'IC et le P\&L perdus à cause de la fenêtre mal alignée, et le taux de détection du test de parité pour chaque
écart injecté.
\item Quel écart est le plus dangereux, et pourquoi~?
\item Comment déploie-t-on une modification de la définition d'une variable~?
\item Que fait-on le jour où le test de parité se déclenche~?
\item Comment intègre-t-on les variables d'un fournisseur dans le magasin~?
\item Que doit le magasin au validateur du modèle (chapitre~21)~?
\item Quel est le coût d'un \omterm{def:ml:data-and-feature-stores:store}{feature store} pour une petite équipe, et quand vaut-il la peine~?
\item En une phrase~: que garantit un \omterm{def:ml:data-and-feature-stores:store}{feature store}~?
\end{enumerate}
\end{problem}

\section{Questions d'entretien}

\begin{interviewq}[$\star$ \iqroles{ingénieur ML}]\label{iq:ml:data-and-feature-stores:1}
Qu'est-ce que l'\omterm{def:ml:data-and-feature-stores:skew}{écart entraînement--service}~? Donnez trois causes.
\end{interviewq}

\begin{interviewq}[$\star\star$ \iqroles{ingénieur ML, développeur}]\label{iq:ml:data-and-feature-stores:2}
Concevez un \omterm{def:ml:data-and-feature-stores:store}{feature store} pour des variables de carnet d'ordres utilisées par un modèle qui passe des ordres à chaque message.
\end{interviewq}

\begin{interviewq}[$\star\star$ \iqroles{chercheur, ingénieur ML}]\label{iq:ml:data-and-feature-stores:3}
L'IC en direct de votre modèle ne vaut que la moitié de celui du backtest. Comment déterminer si les variables sont en cause~?
\end{interviewq}

\begin{interviewq}[$\star\star$ \iqroles{ingénieur ML}]\label{iq:ml:data-and-feature-stores:4}
Qu'est-ce qu'une jointure point-in-time, et comment un \omterm{def:ml:data-and-feature-stores:store}{feature store} la met-il en œuvre~?
\end{interviewq}

\begin{interviewq}[$\star\star$ \iqroles{ingénieur ML, développeur}]\label{iq:ml:data-and-feature-stores:5}
Comment surveillez-vous la \omterm{def:ml:data-and-feature-stores:fresh}{fraîcheur des variables}, et que faites-vous lorsqu'une variable devient obsolète~?
\end{interviewq}

\begin{interviewq}[$\star\star\star$ \iqroles{ingénieur ML}]\label{iq:ml:data-and-feature-stores:6}
Comment feriez-vous coïncider bit à bit une implémentation en ligne et une implémentation hors ligne d'une somme sur fenêtre~?
\end{interviewq}
